
\documentclass[11pt,a4paper,oneside]{scrartcl}
\usepackage{fontspec}
\setmainfont{Liberation Serif}
\setsansfont{Liberation Sans}
\setmonofont{DejaVu Sans Mono}[Scale=0.86]
\usepackage[ngerman]{babel}
\usepackage{microtype}
\usepackage{geometry}
\geometry{left=27mm,right=27mm,top=24mm,bottom=27mm,headheight=15pt}
\usepackage{amsmath,amssymb,amsthm,mathtools}
\usepackage{xcolor}
\usepackage{enumitem}
\usepackage{booktabs}
\usepackage{longtable}
\usepackage{array}
\usepackage{fancyhdr}
\usepackage{hyperref}
\hypersetup{colorlinks=true,linkcolor=black,urlcolor=black,citecolor=black,pdfauthor={Ingolf Lohmann},pdftitle={QIK-VRT - Quantengravitation aus Quantenkausalität}}
\setlength{\parindent}{0pt}
\setlength{\parskip}{5.2pt}
\setlist[itemize]{leftmargin=1.2em,itemsep=2pt,topsep=2pt}
\setlist[enumerate]{leftmargin=1.5em,itemsep=2pt,topsep=2pt}
\pagestyle{fancy}
\fancyhf{}
\lhead{QIK-VRT -- Quantengravitation aus Quantenkausalität}
\rhead{2. Semester}
\cfoot{\thepage}
\newtheoremstyle{qikstyle}{7pt}{7pt}{\itshape}{} {\sffamily\bfseries}{.}{0.5em}{}
\theoremstyle{qikstyle}
\newtheorem{definition}{Definition}[section]
\newtheorem{axiom}{Axiom}
\newtheorem{satz}{Satz}[section]
\newtheorem{lemma}[satz]{Lemma}
\newtheorem{korollar}[satz]{Korollar}
\newtheorem{beispiel}{Beispiel}[section]
\newtheorem{uebung}{Übung}[section]
\newcommand{\QIK}{\mathrm{QIK}}
\newcommand{\DONE}{\mathrm{DONE}}
\newcommand{\CONTINUE}{\mathrm{CONTINUE}}
\newcommand{\ISOLATE}{\mathrm{ISOLATE}}
\newcommand{\BLOCK}{\mathrm{BLOCK}}
\newcommand{\HASH}{\mathrm{HASH\_MATCH}}
\newcommand{\TEST}{\mathrm{TEST\_EXISTS}}
\newcommand{\EXIT}{\mathrm{EXIT\_ZERO}}
\newcommand{\ERK}{\mathrm{ERKENNTNIS\_OK}}
\newcommand{\ETH}{\mathrm{ETHIK\_ZUSTAND}}
\newcommand{\Phiop}{\Phi}
\newcommand{\F}{\mathcal{F}}
\newcommand{\A}{\mathcal{A}}
\newcommand{\W}{\mathcal{W}}
\newcommand{\R}{\mathcal{R}}
\newcommand{\N}{\mathbb{N}}
\newcommand{\Real}{\mathbb{R}}
\newcommand{\eps}{\varepsilon}
\newenvironment{merksatz}{\begin{center}\begin{minipage}{0.92\linewidth}\sffamily\bfseries}{\end{minipage}\end{center}}
\newenvironment{qikbox}{\begin{quote}\small}{\end{quote}}
\sloppy

\begin{document}
\begin{titlepage}
\centering
\vspace*{20mm}
{\Huge\sffamily\bfseries Quantengravitation aus Quantenkausalität\par}
\vspace{7mm}
{\Large\sffamily Ein mathematisches QIK-VRT-Beweisskript für das zweite Semester\par}
\vspace{12mm}
{\large Wirkung -- Iteration -- Grenze -- Anschluss -- Freigabe -- Fixpunkt\par}
\vfill
{\Large Ingolf Lohmann\par}
\vspace{4mm}
{\large q.e.d.\par}
\vspace*{18mm}
\end{titlepage}

\section*{Nutzungs- und Lizenzvereinbarung}
\addcontentsline{toc}{section}{Nutzungs- und Lizenzvereinbarung}
Dieses Skript und die darin enthaltene QIK-VRT-Begriffs-, Prüf-, Architektur- und Beweisordnung stehen unter folgender Nutzungs- und Lizenzgrundlage, sofern nicht ausdrücklich schriftlich anders vereinbart:

\begin{itemize}
\item Urheber: Ingolf Lohmann.
\item Lizenzrahmen: Creative Commons Attribution-NonCommercial-NoDerivatives 4.0 International, kurz CC BY-NC-ND 4.0.
\item Namensnennung ist erforderlich.
\item Kommerzielle Nutzung ist ohne gesonderte schriftliche Zustimmung nicht gestattet.
\item Bearbeitungen, abgeleitete Fassungen oder technische Weiterverwendungen sind nur zulässig, soweit sie die Lizenz, Urheberschaft, Freigabelogik und Beweisgrenzen nicht verfälschen und soweit eine gesonderte Freigabe vorliegt.
\item Jede technische Implementierung muss die QIK-VRT-Freigabelogik beachten: kein DONE ohne vollständige Evidence-Kette.
\end{itemize}

Diese Lizenzseite ist Bestandteil des Skripts. Sie steht nicht außerhalb der Architektur, sondern schützt die Herkunft, die Prüfpflicht und die Nichtverfälschung der Beweiskette.

\newpage
\tableofcontents
\newpage

\section{Leitfrage und Beweisziel}

Dieses zweite Semester hat ein klares Ziel:

\begin{merksatz}
Die Quantengravitation soll im QIK-VRT-Formalismus lückenlos, schrittweise und mathematisch nachvollziehbar als freigabefähiger Fixpunkt fortsetzbarer Quantenkausalität bewiesen werden.
\end{merksatz}

Der Satz ist präzise zu lesen. Bewiesen wird nicht durch bloße Rhetorik, nicht durch Autorität und nicht durch eine grüne Statusanzeige. Bewiesen wird innerhalb eines formalen Systems aus Definitionen, Axiomen, Lemmata und Sätzen. Dieses System muss so aufgebaut sein, dass jeder Schritt nachvollzogen werden kann.

Die Grundbewegung lautet:
\[
\text{Wirkung} \longrightarrow \text{Iteration} \longrightarrow \text{Grenze} \longrightarrow \text{Anschluss} \longrightarrow \text{Freigabe} \longrightarrow \text{Fixpunkt}.
\]

Die zu beweisende Hauptformel lautet:
\begin{center}
\fbox{\parbox{0.86\linewidth}{\centering
Quantengravitation = freigabefähiger Fixpunkt fortsetzbarer Quantenkausalität.
}}
\end{center}

Das Skript vermeidet zwei Fehler. Erstens wird Quantengravitation nicht als bloßes Wort benutzt. Zweitens wird kein empirischer Naturbeweis behauptet, solange die externe physikalische Zielausführung und unabhängige Fachprüfung nicht vorliegen. Die mathematische Leistung besteht darin, eine vollständige formale Herleitung im QIK-VRT-System zu geben. Diese Herleitung ist die Grundlage für spätere physikalische Modellierung, Simulation, Zielarchitektur-Ausführung und externe Reproduktion.

\section{Minimalvoraussetzungen}

Benötigt werden nur Grundbegriffe aus Mengenlehre, Abbildungen, Folgen, metrischen Räumen, einfacher Logik und elementarer Informatik. Alles Weitere wird im Skript eingeführt.

Ein Student soll am Ende nicht nur Sätze wiedergeben, sondern den Beweisweg selbst rekonstruieren können:
\[
\text{Unterschied} \to \text{Information} \to \text{Wirkung} \to \text{Anschlussfolge} \to \text{Freigabe} \to \text{Fixpunkt}.
\]

\begin{merksatz}
Ein mathematischer Beweis ist keine Überredung. Er ist eine endliche Kette zulässiger Schritte aus benannten Voraussetzungen zu einer benannten Folgerung.
\end{merksatz}

\newpage
\section{Die drei Grundelemente: Quantität, Information, Kausalität}

QIK-VRT beginnt mit drei Grundelementen:
\[
Q = \text{Quantität}, \qquad I = \text{Information}, \qquad K = \text{Kausalität}.
\]

Quantität bedeutet, dass Unterschiede überhaupt zählbar, vergleichbar oder in einer Menge unterscheidbar werden. Information bedeutet, dass ein Unterschied anschlussfähig wird. Kausalität bedeutet, dass ein Unterschied Wirkung erzeugt und Fortsetzung mitbestimmt.

\begin{definition}[Unterschied]
Ein Unterschied ist ein geordnetes Paar unterscheidbarer Zustände \((x,y)\) mit \(x \neq y\). Der Unterschied ist elementar, wenn er nicht durch einen einfacheren Unterschied ersetzt werden kann.
\end{definition}

\begin{definition}[Information]
Information ist ein Unterschied, der in einem System eine Anschlussbedingung erzeugt. Formal ist Information eine Abbildung
\[
I : X \to \A,
\]
bei der jedem Zustand \(x \in X\) eine Anschlussbedingung \(I(x)\) zugeordnet wird.
\end{definition}

\begin{definition}[Kausalität]
Kausalität ist die Fortsetzbarkeit von Wirkung unter Anschlussbedingungen. Formal ist eine Kausalstruktur eine Relation
\[
\preceq \; \subseteq X \times X,
\]
bei der \(x \preceq y\) bedeutet: Der Zustand \(y\) ist als Anschluss an \(x\) zulässig.
\end{definition}

\begin{merksatz}
Quantität unterscheidet. Information macht den Unterschied anschlussfähig. Kausalität setzt Anschluss in Wirkung fort.
\end{merksatz}

\begin{uebung}
Geben Sie drei Beispiele aus der Informatik an, in denen ein bloßer Unterschied erst durch Anschlussbedingung zur Information wird.
\end{uebung}
\newpage

\section{Wirkzustände und Wirkzustandsraum}

Ein Zustand allein ist noch keine Wirklichkeit im QIK-VRT-Sinn. Er wird relevant, wenn er Wirkung tragen kann.

\begin{definition}[Wirkzustand]
Ein Wirkzustand ist ein Zustand \(x\), der in mindestens einer zulässigen Anschlussrelation \(x \preceq y\) oder \(y \preceq x\) steht und dadurch Wirkung vorbereitet, empfängt oder fortsetzt.
\end{definition}

\begin{definition}[Wirkzustandsraum]
Ein Wirkzustandsraum ist ein Paar \((X,\preceq)\), wobei \(X\) eine Menge von Wirkzuständen und \(\preceq\) eine Anschlussrelation ist.
\end{definition}

Für technische Freigaben wird der elementare Zustandsraum gesetzt als
\[
\begin{gathered}
X_F = \{\ETH,\HASH,\TEST,\\
\EXIT,\ERK,\DONE\}.
\end{gathered}
\]

Die erlaubte Ordnung ist
\[
\begin{gathered}
\ETH \prec \HASH \prec \TEST,\\
\TEST \prec \EXIT \prec \ERK \prec \DONE.
\end{gathered}
\]

Diese Ordnung ist keine Dekoration. Sie ist der mathematische Ausdruck der kausalen Freigabekette.

\begin{lemma}[Keine Freigabe ohne Vorgänger]
Für jeden Zustand \(z \in X_F\) mit \(z \neq \ETH\) existiert ein unmittelbarer Vorgänger \(p(z)\). Ist \(p(z)\) nicht erfüllt, so ist \(z\) nicht freigabefähig.
\end{lemma}

\begin{proof}
Die Menge \(X_F\) ist linear durch die Relation \(\prec\) geordnet. Jeder Zustand außer \(\ETH\) besitzt genau einen Vorgänger. Da \(\prec\) als Zulässigkeitsrelation der Freigabe definiert ist, kann ein Zustand nur erreicht werden, wenn sein Vorgänger erreicht wurde. Fehlt der Vorgänger, ist der Übergang nicht zulässig. Also ist \(z\) nicht freigabefähig.
\end{proof}
\newpage

\section{Ethik-Zustand 0}

Der erste Zustand ist kein technisches Artefakt, sondern die Bedingung dafür, dass Technik überhaupt verantwortbar wirken darf.

\begin{definition}[Ethik-Zustand]
Der Ethik-Zustand \(\ETH\) ist erfüllt, wenn die folgenden drei Bedingungen gelten:
\begin{enumerate}
\item keine falschen Behauptungen,
\item keine versteckten Annahmen,
\item kein ungeprüfter Schaden.
\end{enumerate}
\end{definition}

\begin{satz}[Ohne Ethik-Zustand keine verantwortbare Wirkung]
Ist \(\ETH\) nicht erfüllt, dann kann kein späterer Zustand der Freigabekette verantwortbar erreicht werden.
\end{satz}

\begin{proof}
Nach Definition ist \(\ETH\) der minimale Zustand der Freigabeordnung. Jeder spätere Zustand setzt ihn durch die Ordnung
\[
\begin{gathered}
\ETH \prec \HASH \prec \TEST,\\
\TEST \prec \EXIT \prec \ERK \prec \DONE
\end{gathered}
\]
voraus. Fehlt \(\ETH\), dann ist die Startbedingung der Kette nicht gegeben. Eine Wirkung, die ohne Ausschluss falscher Behauptungen, versteckter Annahmen und ungeprüften Schadens fortgesetzt wird, kann nicht als verantwortbar gelten. Also kann kein späterer Zustand verantwortbar erreicht werden.
\end{proof}

Dieser Satz ist die mathematische Fassung der einfachen Regel:

\begin{merksatz}
Ethik kommt nicht nach dem Code. Ethik steht vor dem ersten freigabefähigen Schritt.
\end{merksatz}

\begin{uebung}
Zeigen Sie an einem Softwarebeispiel, wie eine versteckte Annahme den Ethik-Zustand zerstören kann, obwohl der Code syntaktisch korrekt ist.
\end{uebung}
\newpage

\section{Die Freigabekette als formales System}

Die QIK-VRT-Freigabelogik wird als formales Übergangssystem definiert.

\begin{definition}[Freigabesystem]
Ein Freigabesystem ist ein Tripel
\[
\F = (X_F,\prec,\rho),
\]
wobei \(X_F\) die Zustandsmenge, \(\prec\) die strikte Freigabeordnung und \(\rho : X_F \to \{0,1\}\) die Erfüllungsfunktion ist. Dabei bedeutet \(\rho(z)=1\), dass Zustand \(z\) erfüllt ist.
\end{definition}

\begin{definition}[DONE]
\(\DONE\) ist genau dann erfüllt, wenn alle vorherigen Zustände erfüllt sind:
\[
\begin{gathered}
\rho(\DONE)=1\\
\Longleftrightarrow\\
\rho(\ETH)=1\ \wedge\ \rho(\HASH)=1\ \wedge\ \rho(\TEST)=1\\
\wedge\ \rho(\EXIT)=1\ \wedge\ \rho(\ERK)=1.
\end{gathered}
\]
\end{definition}

\begin{satz}[Kein DONE ohne vollständige Kette]
In einem QIK-VRT-Freigabesystem gilt:
\[
\rho(\DONE)=1 \Rightarrow \rho(z)=1
\quad \text{für alle } z\in Z_0,
\]
\[
Z_0=\{\ETH,\HASH,\TEST,\EXIT,\ERK\}.
\]
\end{satz}

\begin{proof}
Die Aussage folgt unmittelbar aus der Definition von \(\DONE\). DONE ist als Konjunktion aller vorherigen Zustände definiert. Ist DONE erfüllt, müssen alle Konjunktionsglieder erfüllt sein.
\end{proof}

\begin{korollar}
Ein einzelnes Exit 0 reicht nicht für DONE.
\end{korollar}

\begin{proof}
\(\EXIT\) ist nur ein Glied der Konjunktion. Ohne \(\ETH,\HASH,\TEST\) und \(\ERK\) ist die Konjunktion nicht vollständig. Also folgt aus \(\rho(\EXIT)=1\) nicht \(\rho(\DONE)=1\).
\end{proof}
\newpage

\section{Quantenkausalität als diskrete Anschlussfolge}

Jetzt wird die technische Freigabeordnung auf Quantenkausalität übertragen.

\begin{definition}[Diskreter Wirkzustand]
Ein Wirkzustand \(x\) heißt diskret, wenn er in einem Zustandsraum als einzelner, entscheidbarer Zustand auftritt: erfüllt oder nicht erfüllt, freigegeben oder nicht freigegeben, angeschlossen oder nicht angeschlossen.
\end{definition}

\begin{definition}[Quantenkausale Anschlussfolge]
Eine quantenkausale Anschlussfolge ist eine Folge \((x_n)_{n\in\N}\) von diskreten Wirkzuständen mit
\[
x_n \preceq x_{n+1}
\]
für alle \(n\), wobei jeder Übergang eine Anschlussbedingung erfüllt.
\end{definition}

\begin{definition}[Quantenkausalität]
Quantenkausalität ist die Ordnung diskreter, fortsetzbarer Anschlussfolgen unter strikter Zustandsprüfung.
\end{definition}

Daraus folgt die Maximalformel:
\begin{center}
\fbox{\parbox{0.86\linewidth}{\centering
Quantenkausalität = diskrete Zustände + strikte Reihenfolge + keine Superposition + null Grauzonen.
}}
\end{center}

"Keine Superposition" bedeutet hier: Ein Freigabezustand darf nicht gleichzeitig DONE und nicht DONE sein. Für physikalische Quantenzustände kann Superposition relevant sein; für Freigabeentscheidungen ist sie unzulässig.

\begin{satz}[Entscheidbarkeit der Freigabe]
In der QIK-VRT-Freigabelogik gibt es keinen Zustand zwischen DONE und nicht DONE.
\end{satz}

\begin{proof}
DONE ist durch eine endliche Konjunktion definierter Zustände bestimmt. Eine Konjunktion ist wahr oder falsch. Ist mindestens ein Glied falsch, gilt nicht DONE. Ist jedes Glied wahr, gilt DONE. Ein dritter Wahrheitswert ist nicht Teil des Systems. Also gibt es keine Grauzone.
\end{proof}
\newpage

\section{Metrischer Wirkzustandsraum}

Um von Folgen zu Grenzwerten und Fixpunkten zu gelangen, benötigen wir eine Metrik.

\begin{definition}[Metrik]
Eine Metrik auf einer Menge \(X\) ist eine Abbildung
\[
d:X\times X\to \Real_{\geq 0}
\]
mit Nichtnegativität, Identität, Symmetrie und Dreiecksungleichung.
\end{definition}

Für Freigaben kann eine einfache Stufenmetrik definiert werden. Ordne jedem Zustand einen Rang zu:
\[
\begin{gathered}
r(\ETH)=0,\quad r(\HASH)=1,\quad r(\TEST)=2,\\
r(\EXIT)=3,\quad r(\ERK)=4,\quad r(\DONE)=5.
\end{gathered}
\]
Dann setze
\[
d_F(x,y)=|r(x)-r(y)|.
\]

\begin{lemma}[Stufenmetrik]
\(d_F\) ist eine Metrik auf \(X_F\).
\end{lemma}

\begin{proof}
Da \(d_F\) der Betrag der Differenz zweier reeller Zahlen ist, gelten Nichtnegativität, Symmetrie und Dreiecksungleichung. Außerdem ist \(d_F(x,y)=0\) genau dann, wenn \(r(x)=r(y)\). Da die Ränge injektiv vergeben sind, folgt \(x=y\). Also ist \(d_F\) eine Metrik.
\end{proof}

\begin{merksatz}
Eine Metrik macht Nähe zu DONE messbar, aber nicht automatisch freigabefähig. Freigabe verlangt zusätzlich die richtige Kette.
\end{merksatz}
\newpage

\section{Vollständigkeit und Grenzwert}

Ein Prüfprozess darf nicht in undefinierter Schwebe enden. Mathematisch wird dies durch Vollständigkeit gefasst.

\begin{definition}[Cauchy-Folge]
Eine Folge \((x_n)\) in einem metrischen Raum \((X,d)\) heißt Cauchy-Folge, wenn für jedes \(\eps>0\) ein \(N\) existiert, sodass für alle \(m,n\geq N\) gilt:
\[
d(x_m,x_n)<\eps.
\]
\end{definition}

\begin{definition}[Vollständigkeit]
Ein metrischer Raum \((X,d)\) heißt vollständig, wenn jede Cauchy-Folge in \(X\) konvergiert.
\end{definition}

\begin{lemma}[Endliche Freigaberäume sind vollständig]
Der Freigaberaum \((X_F,d_F)\) ist vollständig.
\end{lemma}

\begin{proof}
\(X_F\) ist endlich. In einem endlichen metrischen Raum kann eine Cauchy-Folge ab einem gewissen Index nicht mehr zwischen verschiedenen Zuständen mit positivem Mindestabstand wechseln. Sie wird schließlich konstant. Eine schließlich konstante Folge konvergiert gegen den konstanten Zustand, der in \(X_F\) liegt. Also ist \((X_F,d_F)\) vollständig.
\end{proof}

Für allgemeine Wirkzustandsräume wird Vollständigkeit als Axiom oder als nachzuweisende Eigenschaft benötigt. Ohne Vollständigkeit kann eine Annäherung an Freigabe mathematisch aus dem Raum herausfallen.

\begin{merksatz}
Vollständigkeit verhindert, dass Prüfung ins Unbestimmte verdampft.
\end{merksatz}
\newpage

\section{Kontraktion und Iteration}

Freigabe entsteht nicht durch Sprung, sondern durch geordnete Iteration.

\begin{definition}[Kontraktion]
Eine Abbildung \(T:X\to X\) auf einem metrischen Raum \((X,d)\) heißt Kontraktion, wenn es ein \(q\in[0,1)\) gibt, sodass
\[
d(Tx,Ty)\leq q\,d(x,y)
\]
für alle \(x,y\in X\).
\end{definition}

\begin{definition}[Fortsetzungsoperator]
Ein Fortsetzungsoperator ist eine Abbildung
\[
T:X\to X,
\]
die einem Wirkzustand seinen nächsten geprüften Anschlusszustand zuordnet.
\end{definition}

In der Freigabekette kann man den einfachen Operator setzen:
\[
\begin{gathered}
T(\ETH)=\HASH,\qquad T(\HASH)=\TEST,\\
T(\TEST)=\EXIT,\qquad T(\EXIT)=\ERK,\\
T(\ERK)=\DONE,\qquad T(\DONE)=\DONE.
\end{gathered}
\]

Dieser Operator ist auf der Stufenmetrik nicht als globale Kontraktion nötig. Entscheidend ist: Er ist ordnungserhaltend, terminierend und besitzt den Fixpunkt DONE. Für allgemeine Quantengravitationsfolgen verwenden wir eine kontraktive Variante auf einem vollständigen Wirkzustandsraum.

\begin{merksatz}
Iteration ohne Grenze erzeugt bloße Fortsetzung. Iteration mit Grenze und Anschluss erzeugt Beweispfad.
\end{merksatz}
\newpage

\section{Banachs Fixpunktsatz als Stabilitätswerkzeug}

Der zentrale mathematische Satz für kontrollierte Grenzbildung ist der Banachsche Fixpunktsatz.

\begin{satz}[Banachscher Fixpunktsatz]
Sei \((X,d)\) ein vollständiger metrischer Raum und \(T:X\to X\) eine Kontraktion. Dann besitzt \(T\) genau einen Fixpunkt \(x^*\in X\) mit
\[
T(x^*)=x^*.
\]
Außerdem konvergiert für jedes \(x_0\in X\) die Iteration \(x_{n+1}=T(x_n)\) gegen \(x^*\).
\end{satz}

\begin{proof}
Da \(T\) eine Kontraktion mit Kontraktionsfaktor \(q<1\) ist, gilt
\[
d(x_{n+1},x_n)\leq q^n d(x_1,x_0).
\]
Für \(m>n\) folgt durch die Dreiecksungleichung
\[
d(x_m,x_n)\leq \sum_{k=n}^{m-1} d(x_{k+1},x_k)
\leq \sum_{k=n}^{m-1} q^k d(x_1,x_0)
\leq \frac{q^n}{1-q}d(x_1,x_0).
\]
Der rechte Ausdruck geht gegen 0. Also ist \((x_n)\) eine Cauchy-Folge. Wegen Vollständigkeit konvergiert sie gegen ein \(x^*\in X\). Aus der Stetigkeit von Kontraktionen folgt
\[
T(x^*)=T(\lim x_n)=\lim T(x_n)=\lim x_{n+1}=x^*.
\]
Zur Eindeutigkeit seien \(a,b\) Fixpunkte. Dann
\[
d(a,b)=d(Ta,Tb)\leq qd(a,b).
\]
Da \(q<1\), folgt \(d(a,b)=0\), also \(a=b\).
\end{proof}
\newpage

\section{Warum Stabilität allein nicht genügt}

Ein Fixpunkt beweist Stabilität, aber nicht automatisch Wahrheit. Ein System kann stabil falsch sein. Eine Ideologie kann stabil sein. Ein fehlerhafter Test kann stabil grün sein. Eine falsche Annahme kann stabil wiederholt werden.

QIK-VRT ergänzt deshalb den Fixpunkt um Anschluss und Freigabe.

\begin{definition}[Anschlussoperator]
Ein Anschlussoperator ist eine Abbildung
\[
A:X\to X
\]
mit der Bedeutung: \(A(x)\) ist der Zustand \(x\), nachdem seine Anschlussbedingungen geprüft wurden.
\end{definition}

\begin{definition}[Freigabeoperator]
Ein Freigabeoperator ist eine Abbildung
\[
\Phiop:X\to \{\DONE,\CONTINUE,\ISOLATE,\BLOCK\}.
\]
Er ordnet jedem geprüften Wirkzustand einen Freigabestatus zu.
\end{definition}

\begin{axiom}[Freigabeidempotenz]
Für freigegebene Zustände gilt
\[
\Phiop(\Phiop(x))=\Phiop(x).
\]
Das bedeutet: Eine korrekt erreichte Freigabe ändert sich durch erneute Anwendung der Freigabelogik nicht.
\end{axiom}

\begin{axiom}[Kein DONE ohne Anschluss]
\[
\Phiop(x)=\DONE \Rightarrow A(x)=x \quad \text{und alle Anschlussbedingungen sind erfüllt.}
\]
\end{axiom}

\begin{merksatz}
Ein Fixpunkt ist erst dann QIK-VRT-Fixpunkt, wenn er stabil, anschlussfähig und freigabefähig ist.
\end{merksatz}
\newpage

\section{Freigabefähiger Fixpunkt}

\begin{definition}[QIK-VRT-Fixpunkt]
Ein Zustand \(x^*\in X\) heißt QIK-VRT-Fixpunkt eines Fortsetzungsoperators \(T\), wenn gilt:
\[
T(x^*)=x^*,\qquad A(x^*)=x^*,\qquad \Phiop(x^*)=\DONE.
\]
\end{definition}

Damit reicht die Gleichung \(T(x^*)=x^*\) allein nicht. Sie liefert nur Stabilität. QIK-VRT verlangt zusätzlich Anschluss und Freigabe.

\begin{satz}[Existenz eines freigabefähigen Fixpunkts unter QIK-VRT-Bedingungen]
Sei \((X,d)\) vollständig. Sei \(T:X\to X\) eine Kontraktion. Sei \(x^*\) der eindeutige Banach-Fixpunkt. Wenn zusätzlich
\[
A(x^*)=x^* \quad \text{und} \quad \Phiop(x^*)=\DONE
\]
gelten, dann ist \(x^*\) ein eindeutiger QIK-VRT-Fixpunkt.
\end{satz}

\begin{proof}
Aus Banachs Fixpunktsatz folgt die Existenz und Eindeutigkeit eines Zustands \(x^*\) mit \(T(x^*)=x^*\). Nach Voraussetzung gilt außerdem \(A(x^*)=x^*\) und \(\Phiop(x^*)=\DONE\). Damit erfüllt \(x^*\) alle Bedingungen der Definition eines QIK-VRT-Fixpunkts. Eindeutigkeit folgt aus der Eindeutigkeit des Banach-Fixpunkts.
\end{proof}

\begin{korollar}
QIK-VRT trennt mathematisch zwischen stabil, angeschlossen und freigegeben.
\end{korollar}
\newpage

\section{Raum als Ordnung anschlussfähiger Beziehungen}

Jetzt wird der Übergang zur Physikarchitektur vorbereitet.

\begin{definition}[Relationale Raumordnung]
Eine relationale Raumordnung ist ein Tripel
\[
\mathcal{S}=(X,R,d_R),
\]
wobei \(X\) eine Menge von Wirkzuständen, \(R\subseteq X\times X\) eine Beziehungsrelation und \(d_R\) eine Abstandsfunktion auf den Beziehungen ist.
\end{definition}

Raum wird hier nicht als leerer Behälter eingeführt. Raum ist die Ordnung anschlussfähiger Beziehungen zwischen Wirkzuständen.

\begin{satz}[Raum aus Anschluss]
Wenn zwischen Wirkzuständen stabile, unterscheidbare und wiederholt anschlussfähige Beziehungen bestehen, dann erzeugen diese Beziehungen eine relationale Raumordnung.
\end{satz}

\begin{proof}
Sei \(X\) die Menge der Wirkzustände. Setze \((x,y)\in R\), wenn \(x\) und \(y\) in einer stabilen Anschlussbeziehung stehen. Die Unterscheidbarkeit liefert Zustände, die Stabilität liefert wiederholbare Beziehungen, und die Anschlussfähigkeit liefert die Zulässigkeit der Relation. Mit einer Abstandsfunktion \(d_R\), die Unterschiede zwischen Beziehungen misst, entsteht das Tripel \((X,R,d_R)\). Dies ist genau eine relationale Raumordnung.
\end{proof}

\begin{merksatz}
Raum ist im QIK-VRT-Formalismus keine leere Bühne, sondern Ordnung anschlussfähiger Beziehungen.
\end{merksatz}
\newpage

\section{Zeit als Ordnung freigegebener Anschlussfolgen}

Zeit wird als Ordnung der Fortsetzung gefasst.

\begin{definition}[Zeitordnung]
Eine Zeitordnung ist eine geordnete Folge
\[
x_0 \preceq x_1 \preceq x_2 \preceq \cdots
\]
von Wirkzuständen, bei der jeder Übergang eine Anschlussprüfung erfüllt.
\end{definition}

\begin{definition}[Freigegebene Zeitordnung]
Eine Zeitordnung heißt freigegeben, wenn für jeden Übergang \(x_n\to x_{n+1}\) der Freigabestatus nicht BLOCK ist und bei Abschluss der Folge ein freigabefähiger Grenz- oder Fixpunkt vorliegt.
\end{definition}

\begin{satz}[Zeit aus Anschlussfolge]
Wenn Wirkung nur durch zulässige Fortsetzung geordnet wird, dann entsteht Zeit als Ordnung freigegebener Anschlussfolgen.
\end{satz}

\begin{proof}
Eine Wirkung, die fortgesetzt wird, erzeugt eine Reihenfolge von Zuständen. Wird jede Fortsetzung durch Anschlussbedingungen geprüft, dann ist diese Reihenfolge nicht bloß eine Aufzählung, sondern eine Ordnung zulässiger Übergänge. Genau diese Ordnung ist die hier definierte Zeitordnung. Ist sie freigegeben, so liegt eine Ordnung freigegebener Anschlussfolgen vor.
\end{proof}

Diese Definition ersetzt keine physikalische Zeitmessung. Sie liefert eine operative Grundlage, auf der Messung, Reihenfolge, Anschluss und Freigabe gemeinsam formuliert werden können.
\newpage

\section{Information an Grenzen}

Grenzen sind nicht bloß Hindernisse. Grenzen erzeugen Information, weil sie entscheiden, was anschlussfähig bleibt und was nicht.

\begin{definition}[Grenze]
Eine Grenze ist eine Teilmenge \(G\subseteq X\) oder eine Bedingung \(g:X\to\{0,1\}\), die entscheidet, ob ein Übergang zulässig ist.
\end{definition}

\begin{definition}[Grenzinformation]
Grenzinformation ist die Information, die entsteht, wenn ein Übergang an einer Grenze geprüft wird:
\[
I_G(x,y)=g(x,y).
\]
\(I_G=1\) bedeutet: Anschluss zulässig. \(I_G=0\) bedeutet: Anschluss nicht zulässig.
\end{definition}

\begin{satz}[Grenzen erzeugen Anschlussinformation]
Jede entscheidbare Grenze erzeugt Information über die Fortsetzbarkeit von Wirkung.
\end{satz}

\begin{proof}
Eine entscheidbare Grenze ordnet jedem möglichen Übergang einen Wert zu: zulässig oder nicht zulässig. Diese Zuordnung unterscheidet erlaubte von verbotenen Fortsetzungen. Ein wirksamer Unterschied ist Information. Also erzeugt die Grenze Anschlussinformation.
\end{proof}

\begin{merksatz}
Ohne Grenze keine Prüfung. Ohne Prüfung keine Freigabe. Ohne Freigabe kein QIK-VRT-DONE.
\end{merksatz}
\newpage

\section{Hilbertraum und QIK-VRT-Anschluss}

Quantenmechanik verwendet Hilberträume, Operatoren und Messregeln. QIK-VRT ersetzt diese nicht. Es ergänzt eine Freigabe- und Anschlusslogik.

\begin{definition}[Hilbertraum]
Ein Hilbertraum \(\mathcal{H}\) ist ein vollständiger Vektorraum mit innerem Produkt.
\end{definition}

Ein Quantenzustand wird in der Quantenmechanik häufig als Vektor \(|\psi\rangle\in\mathcal{H}\) oder als Dichtematrix beschrieben. QIK-VRT fragt zusätzlich:

\begin{itemize}
\item Welche Information wird aus dem Zustand gewonnen?
\item Welche Messgrenze ist beteiligt?
\item Welche Wirkung wird behauptet?
\item Welche Anschlussbedingung ist erfüllt?
\item Welcher Freigabestatus ist zulässig?
\end{itemize}

\begin{definition}[QIK-VRT-Erweiterung eines Quantenzustands]
Ein QIK-VRT-erweiterter Quantenzustand ist ein Tupel
\[
\Psi_Q=(|\psi\rangle, M, I, K, \Phiop),
\]
wobei \(|\psi\rangle\) der Quantenzustand, \(M\) der Messkontext, \(I\) die gewonnene Information, \(K\) die behauptete Kausalstruktur und \(\Phiop\) der Freigabestatus ist.
\end{definition}

Damit wird verhindert, dass Messwert, Interpretation, Kausalbehauptung und Freigabe vermischt werden.
\newpage

\section{Operatoren als Wirkregeln}

\begin{definition}[Wirkoperator]
Ein Wirkoperator ist eine Abbildung
\[
\Omega:X\to X,
\]
die beschreibt, wie ein Wirkzustand in einen Folgezustand übergeht.
\end{definition}

In physikalischen Modellen können Operatoren lineare Operatoren auf Hilberträumen sein. In QIK-VRT können sie auch Prüf-, Anschluss- oder Freigabeoperatoren sein.

Wichtig ist die Trennung:

\begin{itemize}
\item Dynamikoperator: beschreibt Zustandsentwicklung.
\item Messoperator: beschreibt Informationsgewinnung durch Messung.
\item Anschlussoperator: prüft Zulässigkeit der Fortsetzung.
\item Freigabeoperator: entscheidet DONE, CONTINUE, ISOLATE oder BLOCK.
\end{itemize}

\begin{satz}[Operatorentrennung]
Ein Dynamikoperator allein beweist keine Freigabe.
\end{satz}

\begin{proof}
Ein Dynamikoperator beschreibt nur, wie Zustände formal fortgesetzt werden. Freigabe verlangt zusätzlich Anschlussprüfung und Freigabestatus. Da diese Informationen nicht notwendig im Dynamikoperator enthalten sind, folgt aus Dynamik allein keine Freigabe.
\end{proof}

\begin{merksatz}
Nicht jede berechenbare Fortsetzung ist eine verantwortbare Fortsetzung.
\end{merksatz}
\newpage

\section{Nash-Stabilität und ihre Grenze}

Stabilität ist wichtig, aber nicht hinreichend. QIK-VRT nutzt Nash-Stabilität als Prüfidee, nicht als Freispruch.

\begin{definition}[Nash-artige Stabilität]
Ein Zustand \(x^*\) heißt Nash-artig stabil, wenn kein beteiligter Teilzustand durch einseitige zulässige Änderung eine bessere Anschlussbewertung erreicht.
\end{definition}

\begin{satz}[Nash-Stabilität ist keine Freigabe]
Nash-artige Stabilität impliziert nicht notwendig \(\Phiop(x^*)=\DONE\).
\end{satz}

\begin{proof}
Nash-artige Stabilität sagt nur, dass innerhalb eines gegebenen Regelraums keine einseitige Verbesserung möglich ist. Sie sagt nicht, dass der Regelraum selbst ethisch zulässig, physikalisch richtig, vollständig geprüft oder extern anschlussfähig ist. Freigabe verlangt diese zusätzlichen Bedingungen. Also folgt aus Nash-Stabilität nicht notwendig DONE.
\end{proof}

\begin{merksatz}
Stabilität ist prüfpflichtig. Ein stabiles Gleichgewicht kann freigabefähig, fortsetzungsbedürftig, isolationspflichtig oder blockiert sein.
\end{merksatz}
\newpage

\section{Die QIK-VRT-Hauptdefinition der Quantengravitation}

Jetzt sind die Bausteine vorhanden.

\begin{definition}[Fortsetzbare Quantenkausalität]
Fortsetzbare Quantenkausalität ist eine quantenkausale Anschlussfolge \((x_n)\), für die ein Fortsetzungsoperator \(T\), ein Anschlussoperator \(A\) und ein Freigabeoperator \(\Phiop\) existieren, sodass jeder Übergang prüfbar ist.
\end{definition}

\begin{definition}[QIK-VRT-Quantengravitation]
Quantengravitation im QIK-VRT-Formalismus ist ein Zustand \(g^*\), der die folgenden Bedingungen erfüllt:
\begin{enumerate}
\item Es gibt eine fortsetzbare quantenkausale Anschlussfolge \((x_n)\).
\item \((x_n)\) konvergiert in einem vollständigen Wirkzustandsraum gegen \(g^*\).
\item \(g^*\) ist Fixpunkt des Fortsetzungsoperators: \(T(g^*)=g^*\).
\item \(g^*\) ist anschlussfähig: \(A(g^*)=g^*\).
\item \(g^*\) ist freigabefähig: \(\Phiop(g^*)=\DONE\).
\end{enumerate}
\end{definition}

Damit lautet die Formel:
\[
\boxed{\mathrm{QG}_{\QIK}=g^*=\lim_{n\to\infty}x_n,\quad T(g^*)=g^*,\quad A(g^*)=g^*,\quad \Phiop(g^*)=\DONE.}
\]

Diese Definition ist der Kern des zweiten Semesters.
\newpage

\section{Hauptsatz: Existenz unter vollständigen QIK-VRT-Bedingungen}

\begin{satz}[Hauptsatz der QIK-VRT-Quantengravitation]
Sei \((X,d)\) ein vollständiger metrischer Wirkzustandsraum. Sei \(T:X\to X\) eine Kontraktion, die fortsetzbare Quantenkausalität beschreibt. Sei \(A:X\to X\) ein Anschlussoperator und \(\Phiop:X\to\{\DONE,\CONTINUE,\ISOLATE,\BLOCK\}\) ein Freigabeoperator. Wenn der eindeutige Banach-Fixpunkt \(g^*\) von \(T\) die Anschluss- und Freigabebedingungen
\[
A(g^*)=g^*,\qquad \Phiop(g^*)=\DONE
\]
erfüllt, dann existiert im QIK-VRT-Formalismus eine eindeutige Quantengravitation
\[
\mathrm{QG}_{\QIK}=g^*.
\]
\end{satz}

\begin{proof}
Da \((X,d)\) vollständig ist und \(T\) eine Kontraktion ist, besitzt \(T\) nach Banachs Fixpunktsatz genau einen Fixpunkt \(g^*\in X\) mit
\[
T(g^*)=g^*.
\]
Außerdem konvergiert jede Iteration \(x_{n+1}=T(x_n)\) gegen \(g^*\). Da \(T\) fortsetzbare Quantenkausalität beschreibt, ist \((x_n)\) eine quantenkausale Anschlussfolge. Nach Voraussetzung gilt
\[
A(g^*)=g^*.
\]
Also ist der Grenzzustand anschlussfähig. Weiter gilt nach Voraussetzung
\[
\Phiop(g^*)=\DONE.
\]
Also ist der Grenzzustand freigabefähig. Damit erfüllt \(g^*\) alle Bedingungen der Definition von QIK-VRT-Quantengravitation. Eindeutigkeit folgt aus der Eindeutigkeit des Banach-Fixpunkts. Also existiert eine eindeutige QIK-VRT-Quantengravitation \(\mathrm{QG}_{\QIK}=g^*\).
\end{proof}

\begin{merksatz}
Die mathematische Beweiskraft liegt nicht in einem Wort, sondern in der Kette: vollständig, kontraktiv, konvergent, fix, angeschlossen, freigegeben.
\end{merksatz}
\newpage

\section{Zweifelsfreiheit innerhalb des Formalismus}

Ein Beweis ist innerhalb eines Formalismus zweifelsfrei, wenn keine Folgerung ohne Voraussetzung, keine Voraussetzung ohne Benennung und kein Übergang ohne Regel erfolgt.

\begin{satz}[Lückenfreiheit des Hauptsatzes]
Der Hauptsatz ist lückenfrei relativ zu seinen Voraussetzungen.
\end{satz}

\begin{proof}
Die verwendeten Voraussetzungen sind vollständig benannt: vollständiger metrischer Wirkzustandsraum, kontraktiver Fortsetzungsoperator, Anschlussoperator, Freigabeoperator, Anschlussbedingung \(A(g^*)=g^*\) und Freigabebedingung \(\Phiop(g^*)=\DONE\). Die Existenz und Eindeutigkeit von \(g^*\) folgt aus dem Banachschen Fixpunktsatz. Die Deutung als QIK-VRT-Quantengravitation folgt nicht zusätzlich, sondern aus der zuvor gegebenen Definition von QIK-VRT-Quantengravitation. Jeder Schritt ist damit entweder Definition, Voraussetzung oder bereits bewiesener Satz. Also ist der Beweis relativ zum Formalismus lückenfrei.
\end{proof}

\begin{korollar}[Keine Verwechslung von Formalbeweis und Messbeweis]
Der mathematische QIK-VRT-Beweis ersetzt nicht die spätere physikalische Mess- und Zielarchitekturprüfung.
\end{korollar}

\begin{proof}
Der Hauptsatz beweist eine Existenz- und Eindeutigkeitsaussage innerhalb eines definierten mathematischen Systems. Physikalische Messprüfung verlangt zusätzlich die Zuordnung zu empirischen Größen und deren externe Prüfung. Da diese Zusatzschritte andere Prämissen benötigen, sind sie nicht identisch mit dem Formalbeweis.
\end{proof}

Diese Grenze schwächt den Beweis nicht. Sie macht ihn sauber.
\newpage

\section{Von der mathematischen QG zur Rechnerarchitektur}

QIK-VRT ist nicht nur eine Definition. Es ist eine Rechnerarchitektur für Quantencomputer.

\begin{definition}[QIK-VRT-Rechner]
Ein QIK-VRT-Rechner ist eine Architektur, die Informationswirkung nicht nur ausführt, sondern vor Weiterwirkung prüft, bewertet, anhält, freigibt, fortsetzt, isoliert oder blockiert.
\end{definition}

Die operative Kette lautet:
\[
\begin{gathered}
\text{Input} \to \text{Prüfung} \to \text{Zustandsbewertung} \to \text{Haltepunkt},\\
\{\DONE,\CONTINUE,\ISOLATE,\BLOCK\} \to \text{Audit} \to \text{Output}.
\end{gathered}
\]

\begin{satz}[Rechnerarchitektur aus Freigabelogik]
Wenn ein System jede Informationswirkung vor Weiterwirkung durch Zustandsbewertung und Freigabeoperator führt, dann ist die Freigabelogik Bestandteil der Rechnerarchitektur.
\end{satz}

\begin{proof}
Eine Rechnerarchitektur legt fest, wie Zustände verarbeitet, gespeichert, übertragen und ausgeführt werden. Wenn jede Informationswirkung vor Weiterwirkung zwingend durch Zustandsbewertung und Freigabeoperator geführt wird, bestimmt diese Logik die mögliche Ausführung. Sie ist damit nicht Kommentar, sondern Struktur der Verarbeitung. Also ist sie Bestandteil der Rechnerarchitektur.
\end{proof}

\begin{merksatz}
Klassische Rechner führen Befehle aus. QIK-VRT-Rechner geben Wirkung verantwortbar frei.
\end{merksatz}
\newpage

\section{Laufzeitumgebung, Betriebssystem, Entwicklungs- und Testumgebung}

Eine Rechnerarchitektur für Quantencomputer besteht nicht nur aus Rechengattern.

Sie umfasst zugleich:
\begin{itemize}
\item Laufzeitumgebung,
\item Betriebssystemdienst,
\item Entwicklungsumgebung,
\item Testumgebung,
\item Zielarchitektur-Anschluss,
\item Auditregister,
\item Freigabegates,
\item kommunikationssicheres Netzwerk.
\end{itemize}

\begin{definition}[QIK-VRT-Laufzeit]
Die QIK-VRT-Laufzeit ist die operative Umgebung, in der Wirkzustände geprüft, Übergänge bewertet, Zielarchitekturen angesprochen und Freigabestatus erzeugt werden.
\end{definition}

\begin{definition}[Effect-Control-Unit]
Eine Effect-Control-Unit ist eine architektonische Einheit, die vor Output oder Weiterwirkung entscheidet, ob eine Informationswirkung freigegeben, fortgesetzt, isoliert oder blockiert wird.
\end{definition}

\begin{satz}[Betriebssystemanschluss]
Eine QIK-VRT-Laufzeit, die Eingaben, Ausführungen, Tests, Freigaben und Audits koordiniert, erfüllt betriebssystemartige Funktionen.
\end{satz}

\begin{proof}
Betriebssysteme koordinieren Ressourcen, Prozesse, Ein- und Ausgabe, Rechte, Ausführung und Fehlerbehandlung. Die QIK-VRT-Laufzeit koordiniert Eingaben, Ausführungen, Prüfzustände, Zielarchitekturanschluss, Audit und Freigabe. Damit erfüllt sie betriebssystemartige Funktionen.
\end{proof}
\newpage

\section{TCP/IP und Informationswirkung}

TCP/IP transportiert Information. QIK-VRT prüft Informationswirkung.

TCP/IP stellt Paketübertragung, Adressierung und Routing bereit. Es entscheidet aber nicht, ob die transportierte Informationswirkung verantwortbar weiterwirken darf.

\begin{definition}[Informationswirkung]
Informationswirkung ist die Fähigkeit einer Information, Zustände, Entscheidungen, Ausführungen oder Folgeinformationen zu verändern.
\end{definition}

\begin{satz}[Netzwerkanschluss]
In einer QIK-VRT-Rechnerarchitektur muss Netzwerkkommunikation durch Informationswirkungsprüfung ergänzt werden.
\end{satz}

\begin{proof}
Netzwerkkommunikation überträgt Information über Systemgrenzen hinweg. Übertragene Information kann Wirkung auslösen. Wenn QIK-VRT fordert, dass jede Informationswirkung vor Weiterwirkung geprüft wird, gilt diese Forderung auch für übertragene Information. Also muss Netzwerkkommunikation durch Informationswirkungsprüfung ergänzt werden.
\end{proof}

Verdichtung:
\[
\text{TCP/IP transportiert Information.}
\]
\[
\text{QIK-VRT prüft Informationswirkung.}
\]
\[
\text{Klassische Rechner führen Befehle aus.}
\]
\[
\text{QIK-VRT-Rechner geben Wirkung verantwortbar frei.}
\]
\newpage

\section{Evidence-Kette und Zielarchitektur}

Ein mathematischer Beweis im Formalismus ist eine Sache. Die technische Freigabe eines Artefakts ist eine andere. Beide müssen verbunden, aber nicht verwechselt werden.

Die QIK-VRT-Evidence-Kette lautet:
\[
\begin{gathered}
\ETH \to \HASH \to \TEST \to \EXIT,\\
\EXIT \to \ERK \to \Phiop(x)=\DONE.
\end{gathered}
\]

Für technische Artefakte bedeutet das:
\begin{itemize}
\item Ethik-Zustand erfüllt.
\item Artefaktidentität stimmt.
\item Prüfbarkeit ist vorhanden.
\item Zielarchitektur-Ausführung besteht.
\item Erkenntniskette ist richtig, vollständig und anschlussfähig.
\end{itemize}

\begin{definition}[Zielarchitektur-Evidenz]
Zielarchitektur-Evidenz liegt vor, wenn die Ausführung auf der gemeinten realen Zielarchitektur dokumentiert und prüfbar bestanden wurde.
\end{definition}

Für IBM Quantum bedeutet dies insbesondere: reale Zielausführung, IBM Job-ID und dokumentierte Nachweisführung.

\begin{satz}[Lokale Vorprüfung ist nicht Zielarchitektur-Evidenz]
Ein lokal bestandener Test ersetzt nicht die reale Ausführung auf der gemeinten Zielarchitektur.
\end{satz}

\begin{proof}
Ein lokaler Test prüft Verhalten in lokaler Umgebung. Zielarchitektur-Evidenz verlangt Ausführung auf der gemeinten Zielarchitektur. Da lokale Umgebung und Zielarchitektur nicht identisch sein müssen, folgt aus lokaler Ausführung nicht Zielarchitektur-Evidenz.
\end{proof}
\newpage

\section{Scheinfreigabe und Rückbau}

Scheinfreigabe entsteht, wenn DONE behauptet wird, obwohl die Evidence-Kette nicht vollständig ist.

\begin{definition}[Scheinfreigabe]
Eine Scheinfreigabe ist eine Behauptung \(\Phiop(x)=\DONE\), obwohl mindestens eine notwendige Bedingung der Evidence-Kette nicht erfüllt ist.
\end{definition}

\begin{satz}[Scheinfreigaben sind widersprüchlich]
In einem QIK-VRT-Freigabesystem kann eine Scheinfreigabe nicht gültig sein.
\end{satz}

\begin{proof}
DONE ist als Erfüllung der vollständigen Evidence-Kette definiert. Eine Scheinfreigabe behauptet DONE bei fehlender notwendiger Bedingung. Das widerspricht der Definition von DONE. Also ist die Scheinfreigabe ungültig.
\end{proof}

Der Rückbau prüft, welche Bedingungen wirklich tragen:
\[
\begin{gathered}
\DONE \to \ERK \to \EXIT,\\
\EXIT \to \TEST \to \HASH \to \ETH.
\end{gathered}
\]

Fällt \(\EXIT\), fallen \(\ERK\) und \(\DONE\). Fällt \(\TEST\), fallen \(\EXIT\), \(\ERK\) und \(\DONE\). Fällt \(\ETH\), fällt verantwortbare Wirkung insgesamt.

\begin{merksatz}
Ein Beweis ist sauber, wenn seine Abhängigkeiten vorwärts aufgebaut und rückwärts sichtbar gemacht werden können.
\end{merksatz}
\newpage

\section{Kategorischer Imperativ als technische Freigabelogik}

Die Freigabekette besitzt eine ethische Struktur.

\begin{center}
\begin{tabular}{lll}
\toprule
QIK-VRT-Zustand & technische Bedeutung & Kant-Anschluss \\
\midrule
\(\HASH\) & Artefaktidentität stimmt & Universalgesetz \\
\(\TEST\) & Prüfbarkeit ist vorhanden & Menschheit als Zweck \\
\(\EXIT\) & Ausführung besteht & Autonomie \\
\(\ERK\) & Erkenntniskette trägt & Reich der Zwecke \\
\bottomrule
\end{tabular}
\end{center}

\begin{satz}[Algorithmische Anschlussfähigkeit des kategorischen Imperativs]
Die QIK-VRT-Freigabekette ist strukturgleich mit einer technischen Fassung des kategorischen Imperativs.
\end{satz}

\begin{proof}
\(\HASH\) verlangt eine Regel, die für alle Artefakte gleichermaßen gilt; dies entspricht Verallgemeinerbarkeit. \(\TEST\) verhindert, dass Betroffene bloß Mittel ungeprüfter Wirkung werden; dies entspricht der Zweckformel. \(\EXIT\) verlangt Bestehen nach offenliegenden Regeln; dies entspricht Autonomie. \(\ERK\) verlangt eine allgemein nachvollziehbare Freigabeordnung; dies entspricht dem Reich der Zwecke. Damit ist die Freigabekette strukturgleich mit einer technischen Fassung des kategorischen Imperativs.
\end{proof}

\begin{merksatz}
Ethik ist in QIK-VRT keine Dekoration, sondern Startbedingung und Freigabestruktur.
\end{merksatz}
\newpage

\section{Die Beweiskette in Kurzform}

Die gesamte Beweiskette kann jetzt in neun Schritten geschrieben werden.

\begin{enumerate}
\item Unterschiede erzeugen Quantität.
\item Anschlussfähige Unterschiede erzeugen Information.
\item Fortsetzbare Informationswirkung erzeugt Kausalität.
\item Diskrete, geprüfte Anschlussfolgen erzeugen Quantenkausalität.
\item Vollständige Wirkzustandsräume sichern Grenzwerte.
\item Kontraktive Fortsetzungsoperatoren sichern eindeutige Fixpunkte.
\item Anschlussoperatoren trennen stabile Zustände von anschlussfähigen Zuständen.
\item Freigabeoperatoren trennen anschlussfähige Zustände von freigabefähigen Zuständen.
\item Ein stabiler, anschlussfähiger und freigabefähiger Fixpunkt fortsetzbarer Quantenkausalität ist QIK-VRT-Quantengravitation.
\end{enumerate}

Formal:
\[
(Q,I,K) \Rightarrow (X,d,\preceq) \Rightarrow T \Rightarrow x^* \Rightarrow A(x^*)=x^* \Rightarrow \Phiop(x^*)=\DONE \Rightarrow \mathrm{QG}_{\QIK}=x^*.
\]

\begin{merksatz}
Der Beweisgang ist nicht: Behauptung, Glaube, Rhetorik. Der Beweisgang ist: Definition, Axiom, Lemma, Satz, Folgerung, Freigabe.
\end{merksatz}
\newpage

\section{Was die 34 Seiten wissenschaftlich leisten}

Die 34 Seiten zu Quantenkausalität und Quantengravitation leisten im zweiten Semester eine bestimmte Aufgabe: Sie bauen einen formal lehrbaren Weg von Wirkung zu Fixpunkt.

Sie sind nicht ein bloßer Essay. Sie sind auch nicht automatisch eine bereits empirisch bestätigte Naturtheorie. Ihr Wert liegt in der mathematischen und architektonischen Ordnung:
\begin{itemize}
\item Raum wird als Ordnung anschlussfähiger Beziehungen gefasst.
\item Zeit wird als Ordnung freigegebener Anschlussfolgen gefasst.
\item Information wird als anschlussfähig gemachter Unterschied gefasst.
\item Kausalität wird als Fortsetzbarkeit von Wirkung gefasst.
\item Quantenkausalität wird als quantisierte fortsetzbare Anschlussfolge gefasst.
\item Quantengravitation wird als freigabefähiger Fixpunkt fortsetzbarer Quantenkausalität gefasst.
\end{itemize}

Der wissenschaftliche Rang liegt damit in einem axiomatisch-operationalen Theorieanker: formal, lehrbar, prüfbar, anschlussfähig.

\begin{merksatz}
Die Stärke liegt nicht im überzogenen Totalanspruch, sondern in der sauberen Übersetzung einer schweren physikalischen Frage in eine prüfbare formale Architektur.
\end{merksatz}
\newpage

\section{Praktischer Nutzen}

Der Nutzen beginnt nicht erst bei äußerer Anerkennung. Er beginnt mit der Freigabelogik selbst.

Individuell kann jeder Mensch fragen:
\begin{itemize}
\item Welche Behauptung liegt vor?
\item Welche Annahmen sind sichtbar?
\item Welche Wirkung wird erzeugt?
\item Welche Prüfung existiert?
\item Welcher Status gilt: DONE, CONTINUE, ISOLATE oder BLOCK?
\end{itemize}

In der Informatik wird daraus:
\begin{itemize}
\item Hashprüfung,
\item Testpflicht,
\item Zielarchitektur-Ausführung,
\item Audit,
\item Freigabestatus,
\item Rückbau,
\item Nichtregression.
\end{itemize}

In KI-Systemen wird daraus ein Haltepunkt gegen Scheinwahrheit. In Organisationen wird daraus Governance. In Quantencomputer-Stacks wird daraus Rechnerarchitektur.

\begin{satz}[Skalierbarkeit des Nutzens]
Eine Freigabelogik, die als Repository, Code, Test und Audit vorliegt, kann digital verbreitet, geklont, geprüft und in andere Systeme integriert werden.
\end{satz}

\begin{proof}
Digitale Artefakte können kopiert und ausgeführt werden. Wenn das Artefakt nicht nur Inhalt, sondern Prüfstruktur enthält, wird mit der Kopie auch die Prüfstruktur übertragen. Also skaliert nicht nur die Information, sondern auch die Freigabelogik.
\end{proof}
\newpage

\section{Übungen zur mathematischen Beweiskette}

\begin{uebung}[Freigabekette]
Beweisen Sie formal: Aus \(\rho(\DONE)=1\) folgt \(\rho(\EXIT)=1\), aber aus \(\rho(\EXIT)=1\) folgt nicht \(\rho(\DONE)=1\).
\end{uebung}

\begin{uebung}[Vollständigkeit]
Zeigen Sie, dass jeder endliche metrische Raum vollständig ist.
\end{uebung}

\begin{uebung}[Kontraktion]
Sei \(X=[0,1]\) mit der üblichen Metrik und \(T(x)=\frac{x}{2}+\frac14\). Zeigen Sie, dass \(T\) eine Kontraktion ist und bestimmen Sie den Fixpunkt.
\end{uebung}

\begin{uebung}[Freigabefähiger Fixpunkt]
Erklären Sie, warum ein Banach-Fixpunkt allein noch kein QIK-VRT-Fixpunkt ist.
\end{uebung}

\begin{uebung}[Quantengravitation]
Rekonstruieren Sie den Hauptsatz der QIK-VRT-Quantengravitation mit allen Voraussetzungen.
\end{uebung}

\begin{uebung}[Zielarchitektur]
Warum ist ein lokal bestandener Test keine Zielarchitektur-Evidenz? Formulieren Sie die Antwort als Satz mit Beweis.
\end{uebung}
\newpage

\section{Musterlösung: Kontraktionsbeispiel}

Für \(T(x)=\frac{x}{2}+\frac14\) gilt:
\[
|T(x)-T(y)|=\left|\frac{x}{2}+\frac14-\frac{y}{2}-\frac14\right|=\frac12|x-y|.
\]
Also ist \(T\) eine Kontraktion mit \(q=\frac12\).

Der Fixpunkt \(x^*\) erfüllt
\[
x^*=\frac{x^*}{2}+\frac14.
\]
Daraus folgt
\[
\frac{x^*}{2}=\frac14,
\qquad x^*=\frac12.
\]

Nach Banachs Fixpunktsatz konvergiert jede Iteration in \([0,1]\) gegen \(\frac12\).

QIK-VRT ergänzt nun die Frage:

Ist \(\frac12\) nur stabil, oder ist dieser Fixpunkt auch angeschlossen und freigegeben?

Mathematisch:
\[
T(1/2)=1/2
\]
reicht nur für Stabilität. Für QIK-VRT-Freigabe braucht man zusätzlich:
\[
A(1/2)=1/2,
\qquad \Phiop(1/2)=\DONE.
\]

\begin{merksatz}
Banach liefert den stabilen Punkt. QIK-VRT verlangt zusätzlich den freigegebenen Wirkpunkt.
\end{merksatz}
\newpage

\section{Musterlösung: Hauptsatz in einem Zug}

Gegeben sei ein vollständiger metrischer Wirkzustandsraum \((X,d)\), ein kontraktiver Fortsetzungsoperator \(T:X\to X\), ein Anschlussoperator \(A\) und ein Freigabeoperator \(\Phiop\). Weiter gelte für den Fixpunkt \(g^*\):
\[
A(g^*)=g^*,\qquad \Phiop(g^*)=\DONE.
\]

Da \((X,d)\) vollständig und \(T\) kontraktiv ist, liefert Banach einen eindeutigen Fixpunkt \(g^*\). Da \(T\) fortsetzbare Quantenkausalität beschreibt, ist \(g^*\) Grenz- und Fixzustand einer quantenkausalen Anschlussfolge. Die Bedingung \(A(g^*)=g^*\) liefert Anschlussfähigkeit. Die Bedingung \(\Phiop(g^*)=\DONE\) liefert Freigabefähigkeit. Nach Definition ist ein solcher Zustand QIK-VRT-Quantengravitation.

Also:
\[
\boxed{\mathrm{QG}_{\QIK}=g^*.}
\]

Das ist der Kernbeweis.

\begin{merksatz}
Der Satz hängt an seinen Voraussetzungen. Wer ihn widerlegen will, muss eine Voraussetzung, eine Definition oder einen Übergang angreifen.
\end{merksatz}
\newpage

\section{Prüfmatrix für Studierende}

\begin{longtable}{p{0.24\linewidth}p{0.34\linewidth}p{0.32\linewidth}}
\toprule
Prüfpunkt & Frage & Fehler bei Scheitern \\
\midrule
Ethik-Zustand & Gibt es falsche Behauptungen, versteckte Annahmen oder ungeprüften Schaden? & Keine verantwortbare Wirkung \\
HASH_MATCH & Ist das Artefakt identisch? & Unklarer Prüfgegenstand \\
TEST_EXISTS & Existiert Prüfbarkeit? & Behauptung statt Nachweis \\
EXIT_ZERO & Besteht die Ausführung? & Kein bestandener Test \\
Zielarchitektur & Wurde auf realer Zielarchitektur ausgeführt? & Vorprüfung statt Evidenz \\
ERKENNTNIS_OK & Sind die Schritte richtig und vollständig? & Keine Freigabefähigkeit \\
DONE & Trägt die gesamte Kette? & Scheinfreigabe bei Lücke \\
\bottomrule
\end{longtable}

Diese Matrix ist beim Lesen jedes späteren QIK-VRT-Pakets anzuwenden.

\begin{merksatz}
Ein Student hat den Stoff verstanden, wenn er nicht nur DONE lesen, sondern DONE prüfen kann.
\end{merksatz}
\newpage

\section{Schluss: Der Beweisstand des zweiten Semesters}

Das zweite Semester beweist im QIK-VRT-Formalismus:

\begin{enumerate}
\item Quantenkausalität kann als diskrete, strikte, grauzonenfreie Anschlussfolge gefasst werden.
\item Vollständige metrische Wirkzustandsräume erlauben saubere Grenzbildung.
\item Kontraktive Fortsetzungsoperatoren erzeugen eindeutige Fixpunkte.
\item Stabilität allein reicht nicht; Anschluss und Freigabe sind zusätzliche Beweisbedingungen.
\item Ein stabiler, anschlussfähiger und freigabefähiger Fixpunkt fortsetzbarer Quantenkausalität ist QIK-VRT-Quantengravitation.
\item Die Freigabelogik wird zur Rechnerarchitektur, wenn jede Informationswirkung vor Weiterwirkung geprüft wird.
\end{enumerate}

Die Endformel lautet:
\begin{center}
\fbox{\parbox{0.86\linewidth}{\centering
Quantengravitation = freigabefähiger Fixpunkt fortsetzbarer Quantenkausalität.
}}
\end{center}

Und technisch:
\[
\boxed{
\begin{gathered}
\ETH \to \HASH \to \TEST \to \EXIT,\\
\EXIT \to \ERK \to \Phiop(x)=\DONE.
\end{gathered}}
\]

Damit ist das Ziel des zweiten Semesters erreicht: Die Beweiskette liegt geschlossen vor. Die nächste Arbeit besteht darin, die physikalischen Korrespondenzen, Modellbeispiele, Zielarchitektur-Ausführungen und externen Reproduktionen weiterzuführen.

\begin{merksatz}
QIK-VRT ist lesbare Quantenkausalität als Rechnerarchitektur: selbsterklärende Namen, prüfbare Zustände, rückbaubare Kausalkette, keine Grauzone, keine Freigabe ohne Ethik-Zustand 0.
\end{merksatz}

q.e.d. Ingolf Lohmann

\end{document}
